\begin{lemma}[Claim 4.2: strong interchangeability survives an earlier Kempe change]
Let \(s,t\) be PE slots with respective supporting vertices
\(v_s,v_t\) and connecting-color pairs
\(S_s=\{\gamma_s,\delta_s\}\) and
\(S_t=\{\gamma_t,\delta_t\}\), with these four colors equal to the
corresponding exact gamma and delta fields of the certified slots.  Let
\(\varphi\) be a proper coloring
leaving exactly \(e_0\) uncolored, and let \(P\) be an arbitrary
\((\gamma_t,\delta_t)\)-path component.  Suppose that both displayed
color pairs are strongly interchangeable at their respective trees.
Assume either that \(v_s\ne v_t\), in which case
\(S_s\cap S_t=\varnothing\), or that \(s<t\).  Then the coloring
\[
       \varphi^*=\varphi/P
\]
is again proper, leaves exactly \(e_0\) uncolored, and
\(\gamma_s,\delta_s\) remain \(T_s\)-strongly interchangeable under
\(\varphi^*\).
\end{lemma}
\begin{proof}
Write \(K=P\) for the exact connected component through the supplied
root.  By \noderef{TwoColorPathComponent},
\(\gamma_t\ne\delta_t\).  Hence
\noderef{TwoColorComponentSwapProper}, applied to the assumed properness
of \(\varphi\), proves that
\(\varphi^*=\noderef{TwoColorComponentSwap}(\varphi,K)\) is proper.
The same definition shows pointwise that a value is \(\mathsf{none}\)
after the switch exactly when it was \(\mathsf{none}\) before the switch:
on an edge node of \(K\) the definition exchanges the two \(\mathsf{some}\)
values, and off \(K\) it changes nothing.  The exact-one-uncolored
hypothesis \noderef{IsOneEdgeUncolored} therefore remains true with the
same edge \(e_0\).

We prove the strong-interchangeability assertion using the same
first-non-RE witness supplied for slot \(s\) by
\noderef{ETTStrongInterchangeable}.  The extension kinds, supporting
vertex, and minimality clauses depend only on the certified series and are
unchanged by a Kempe switch.  Fix
\(x\in V(T_s)\setminus\{v_s\}\), and let \(C_x\) be its
\((\gamma_s,\delta_s)\)-component under \(\varphi\).  The strong
interchangeability hypothesis at \(s\) says that \(C_x\) is a cycle,
avoids \(v_s\), and is contained in \(T_{\epsilon(s)}\).  It remains to
show that this exact component is unchanged by the switch on \(K\).

First suppose \(v_s\ne v_t\).  The explicit cross-class hypothesis says
that the finite pairs \(\{\gamma_s,\delta_s\}\) and
\(\{\gamma_t,\delta_t\}\) are disjoint.  Expanding
\noderef{TwoColorComponentSwap}, only edges whose old color is
\(\gamma_t\) or \(\delta_t\) can change, and their new color is the other
member of that same pair.  Thus the set of edges colored
\(\gamma_s\) or \(\delta_s\), and every incidence between such an edge
and a vertex, is identical under \(\varphi\) and \(\varphi^*\).
By \noderef{TwoColorComponent}, the whole component \(C_x\) is therefore
identical under the two colorings.  Its cycle witness, avoidance of
\(v_s\), and containment in \(T_{\epsilon(s)}\) all persist.

It remains to treat \(v_s=v_t\).  The alternative in the statement then
gives \(s<t\).  Let \(r=\epsilon(s)\).  Slot \(t\) is a PE by the
strong-interchangeability hypothesis at \(t\), so it is not an RE.  If
\(t<r\), the minimality clause for \(r\) would say that slot \(t\) is an
RE, a contradiction.  Hence \(r\leq t\).  Every positive constructor of
\noderef{TashkinovSeriesStep} makes its current tree an ordered initial
segment of its output tree.  Induction from \(r\) through \(t\), using
\noderef{TreeSequenceInitialSegment}, therefore gives
\[
             V(T_r)\subseteq V(T_t).                           \tag{1}
\]
Consequently the cycle \(C_x\), which is contained in \(T_r\), is
contained in \(T_t\); it also avoids the common vertex
\(v_s=v_t\).

We next show that the path component \(K\) meets \(V(T_t)\) in at most
\(v_t\).  If a vertex \(y\ne v_t\) of \(T_t\) lay in \(K\), strong
interchangeability at \(t\) would make the component through \(y\) a
cycle.  Membership of \(y\) in \(K\) identifies that component with
\(K\), contradicting the path-component clause for \(K\).  This uses the
definitions \noderef{TwoColorComponentCycle} and
\noderef{TwoColorPathComponent}: every node on the former's simple closed
walk has two distinct walk neighbors, while the latter has a node with a
subsingleton neighbor set.  Since every vertex of \(C_x\) lies in \(T_t\)
by (1) and \(C_x\) avoids \(v_t\), the two components have no common
vertex node.  They have no common edge node either.  Indeed, an edge node
on the cycle has an endpoint vertex adjacent to it in
\noderef{TwoColorIncidenceGraph}; if that edge node also belonged to
\(K\), connectedness would put the same endpoint vertex in \(K\), contrary
to vertex-node disjointness.  Thus the vertex and edge nodes of \(K\) and
\(C_x\) are disjoint.

No edge of \(C_x\) is switched.  Moreover, no switched edge can acquire a
\(\gamma_s\)- or \(\delta_s\)-incidence at a vertex of \(C_x\): such an
edge has both endpoint nodes in the component \(K\), while \(K\) is
disjoint from \(C_x\).  Therefore the two-color incidence graph has the
same connected component through \(x\) after the switch.  The old cycle
witness from \noderef{TwoColorComponentCycle} still has exactly that
support, still avoids \(v_s\), and still lies in
\(T_{\epsilon(s)}\).

The conclusion holds for every
\(x\in V(T_s)\setminus\{v_s\}\).  Together with the unchanged PE and
first-non-RE clauses, this is precisely
\noderef{ETTStrongInterchangeable} for
\(\gamma_s,\delta_s\) under \(\varphi^*\), completing the proof.
\end{proof}
